\section{Place chunk ends beside encoding ends}\label{utf8-x-chunks}
The fixture divides eleven bytes into four calls of lengths two, three, two, and
four. The byte offset identifies a location in the whole stream, independent of
which chunk currently contains it. Compare the call boundaries above the row with
the scalar encoding spans below it.

\begin{figure}[htbp]
\centering
\begin{tikzpicture}[alt={Eleven hexadecimal bytes 41 e2 82 ac f0 9f 99 82 00 c3 a9 occupy offsets zero through ten. Chunk ends after offsets one, four and six split the euro and four-byte scalar encodings. Chunk sizes are two, three, two and four. Scalar spans are byte zero for U+0041, one through three for U+20AC, four through seven for U+1F642, byte eight for U+0000, and nine through ten for U+00E9.},x=10mm,y=12mm]
\foreach \i/\value in {0/41,1/e2,2/82,3/ac,4/f0,5/9f,6/99,7/82,8/00,9/c3,10/a9} {
 \draw (\i,0) rectangle +(1,1);
 \node[font=\scriptsize\ttfamily] at (\i+0.5,0.5) {\value};
 \node[above,font=\scriptsize] at (\i+0.5,1) {\i};
}
\foreach \cut in {2,5,7} {\draw[dashed] (\cut,1.35)--(\cut,2.1);}
\node[above,font=\small] at (1,1.7) {Call 0};
\node[above,font=\small] at (3.5,1.7) {Call 1};
\node[above,font=\small] at (6,1.7) {Call 2};
\node[above,font=\small] at (9,1.7) {Call 3};
\foreach \left/\right/\name in {0/1/0041,1/4/20AC,4/8/1F642,8/9/0000,9/11/00E9} {
 \draw (\left+0.05,-0.15)--(\left+0.05,-0.3)--(\right-0.05,-0.3)--(\right-0.05,-0.15);
 \node[below,font=\scriptsize] at ({(\left+\right)/2},-0.35) {U+\name};
}
\end{tikzpicture}
\caption{Delivery and encoding boundaries are independent. Cell labels are hexadecimal bytes; positions and call indices are decimal.}
\end{figure}
After the first call, the scanner has accepted one complete scalar and the lead of
a second. After the second, it has completed that second scalar and accepted the
next lead. The third call advances the pending four-byte encoding without completing
it. The fourth finishes that encoding and then accepts U+0000 and U+00E9.
